\clearpage
\phantomsection
\addcontentsline{toc}{chapter}{Tóm tắt}
\chapter*{Tóm tắt}

\noindent\textbf{Tóm tắt:}
Báo cáo nghiên cứu và mô phỏng lại lớp tấn công giải ẩn danh (de-anonymization) trên các tập dữ liệu thưa, dựa trên công trình kinh điển \textit{``Robust De-anonymization of Large Sparse Datasets''} của Narayanan và Shmatikov - cũng chính là cuộc tấn công đã phá vỡ tính ẩn danh của tập dữ liệu Netflix Prize. Chúng em phân tích mô hình hình thức của bài toán (độ thưa, độ tương tự, định nghĩa rò rỉ riêng tư), thuật toán Scoreboard-RH cùng tiêu chí \textit{eccentricity}, sau đó xây dựng một ứng dụng web hoàn chỉnh (FastAPI + React) mô phỏng kẻ tấn công trên tập dữ liệu công khai MovieLens 100k.

Bằng thực nghiệm thật trên 943 người dùng và 100.000 lượt đánh giá, chúng em chứng minh: chỉ với 8 bộ phim + ngày xem, thuật toán định danh đúng 98\% người dùng; ngay cả khi không có thông tin ngày vẫn đạt 93\%; và với thông tin nhiễu ($\pm 1$ điểm, $\pm 14$ ngày) vẫn đạt 72\% - khẳng định tính \textit{bền vững} của tấn công. Khi ưu tiên các bộ phim hiếm, chỉ 2 phim đã đủ định danh 59\% người dùng. Báo cáo cũng so sánh, tối ưu và nâng cấp ứng dụng demo cho đúng công thức trong bài báo gốc, rồi thảo luận các biện pháp phòng vệ (k-anonymity, riêng tư vi phân, hạn chế công bố micro-data).

\vspace{0.5cm}
\noindent\textit{\textbf{Từ khóa:}} \textit{giải ẩn danh, dữ liệu thưa, quyền riêng tư, Netflix Prize, MovieLens, Scoreboard-RH, eccentricity.}
